\documentclass[../../../include/open-logic-section]{subfiles}

\begin{document}

\olfileid{sth}{card-arithmetic}{fix}
\olsection{$\aleph$-স্থিরবিন্দু}

\olref[spine][]{chap}-এ বলেছিলাম, প্রতিস্থাপনের
ন্যায্যতা দরকার: এই স্বতঃসিদ্ধে সেটের স্তরক্রম
বেশ উঁচু হয়। কিছু অঙ্কবাচক পাটিগণিত জানার পরে
স্তরক্রমের উচ্চতা একটু বোঝানো যায়।

স্পষ্টত $0<\aleph_0$, $1<\aleph_1$,
$2<\aleph_2$\ldots; এবং প্রতি ধাপে আকারের
পার্থক্য আরও \emph{বাড়ে}। তাই অনুমান করতে
ইচ্ছে হয়, প্রত্যেক ক্রমসংখ্যা $\kappa$-র জন্য
$\kappa<\aleph_\kappa$।

কিন্তু $\ZFC$-এ এই অনুমান \emph{ভুল}। বরং
\emph{$\aleph$-স্থিরবিন্দু} আছে, অর্থাৎ এমন
অঙ্কবাচক সংখ্যা $\kappa$, যার জন্য
$\kappa=\aleph_\kappa$।

\begin{prop}\ollabel{alephfixed}
একটি $\aleph$-স্থিরবিন্দু আছে।
\end{prop}

\begin{proof}
পুনরাবৃত্তি ব্যবহার করে সংজ্ঞায়িত করি:
\begin{align*}
	\kappa_0 &= 0\\
	\kappa_{n+1} &= \aleph_{\kappa_n}\\
	\kappa&= \bigcup_{n < \omega}\kappa_n
\end{align*}
\olref[cardinals][classing]{unioncardinalscardinal}
অনুযায়ী $\kappa$ একটি অঙ্কবাচক সংখ্যা। এখন:
\[
	\kappa= \bigcup_{n < \omega} \kappa_{n+1} =
	\bigcup_{n < \omega}\aleph_{\kappa_n} =
	\bigcup_{\alpha < \kappa}\aleph_\alpha = \aleph_\kappa
\]
% BN-SRC-813 TARGET-CORRECTION-BEGIN
\par\noindent\textbf{পাঠ-সংশোধন।} আলেফ ও বেথ উভয় অপেক্ষকই সূচকে কঠোরভাবে বর্ধমান এবং
অশূন্য সীমায় আগের মানগুলির সংযোগ; আবেশে তাদের মান সূচকের
চেয়ে ছোট নয়। আলেফের ক্ষেত্রে $\kappa_0<\kappa_1$, আর
কঠোর বৃদ্ধি থেকে $\kappa_n<\kappa_{n+1}$। প্রতিস্থাপন
দিয়ে এই গণনাযোগ্য পদগুলির সেট পাওয়া যায়। তার সংযোগ
$\kappa$ একটি অশূন্য সীমা ক্রমসংখ্যা এবং পদগুলি তাতে সহশেষ।
তাই প্রদর্শিত তৃতীয় সমতা বৈধ। এছাড়া যেকোনো আলেফ-স্থিরবিন্দু
$\rho$-র জন্য আবেশে $\kappa_n\leq\rho$, ফলে $\kappa\leq\rho$:
এটিই ক্ষুদ্রতম আলেফ-স্থিরবিন্দু। বেথের ক্ষেত্রেও একই প্রমাণ চলে।
% BN-SRC-813 TARGET-CORRECTION-END
%	By construction, $\kappa$ is the least cardinal greater than each
%	$\kappa_n$. So $\aleph_\kappa$ is the least cardinal greater than
%	each $\aleph_{\kappa_n}$, and hence greater than each
%	$\kappa_{n+1}$. But equally $\kappa$ is the least cardinal greater
%	than each $\kappa_{n+1} = \aleph_{\kappa_n}$. So $\kappa=
%	\aleph_\kappa$.
\end{proof}

বুলস একবার ঠিক এই নির্মিত $\aleph$-স্থিরবিন্দু
নিয়ে একটি প্রবন্ধ লিখেছিলেন। প্রবন্ধের শুরুতে
$\kappa$-র অস্তিত্বের কথা বলার পর তিনি লেখেন:
\begin{quote}
	[সেটতত্ত্বের সঙ্গে আগে পরিচয় না থাকলে $\kappa$]
	একটি \emph{বেশ বড়} সংখ্যা বলে মনে হবে। আমার
	মনে হয়, এতই বড় যে, অন্তত $\kappa$-সংখ্যক বস্তু
	আছে বলে দাবি করে এমন যেকোনো তত্ত্বের সত্যতা
	নিয়েই প্রশ্ন ওঠে।
	\citep[p.~257]{Boolos2000}
\end{quote}
প্রবন্ধের শেষে তিনি প্রশ্ন করেন:
\begin{quote}
	[আমাদের কি] সন্দেহ হয় যে গল্পের শুরুতে যা-ই
	থাকুক, এত দূর এসে চাকা কেবল ঘুরছে, আর বাস্তবে
	যা ঘটে তার কোনো বিবরণ আমরা শুনছি না?
	\citep[p.~268]{Boolos2000}
\end{quote}
আমরা যদি সত্যিই ``বাস্তবে যা ঘটে'' তার সীমা
ছাড়িয়ে যাই, তবে দোষ সরাসরি প্রতিস্থাপনের।
এই স্বতঃসিদ্ধেই আমাদের প্রমাণটি কাজ করে।
সেক্ষেত্রে বুলসকে কয়েক দশক আগে করা তাঁর
``কোনো অবাঞ্ছিত'' ফল নেই—এই প্রতিস্থাপন-সংক্রান্ত
দাবি আবার ভেবে দেখতে হতো
(দেখো \olref[replacement][extrinsic]{sec})।

কিন্তু $\kappa$-র অস্তিত্ব কি এত খারাপ? এখানে
রাসেলের \emph{ট্রিস্ট্রাম শ্যান্ডির কূটাভাস}
ভাবলে সুবিধা হতে পারে। ট্রিস্ট্রাম শ্যান্ডি
দিনলিপিতে নিজের জীবন লিখে রাখে, কিন্তু একটি
দিনের কথা লিখতে তার এক বছর লাগে। প্রতি বছর
সে আরও পিছিয়ে পড়ে: এক বছর পরে একটি দিনের
কথা লেখা হয়েছে, অথচ ৩৬৪ দিনের কথা লেখা
হয়নি; দুই বছর পরে দুটি দিনের কথা লেখা হয়েছে,
৭২৮ দিনের কথা লেখা হয়নি; তিন বছর পরে
তিনটি দিনের কথা লেখা হয়েছে, ১০৯২ দিনের
কথা লেখা হয়নি\dots\footnote{অধিবর্ষ উপেক্ষা করছি।}
তবু ট্রিস্ট্রাম যদি \emph{অমর} হয়, তবে শেষ পর্যন্ত
প্রত্যেক দিনের কথাই লিখবে: জীবনের $n$-তম বছরের
মধ্যে $n$-তম দিনের কথা লিখবে। ফলে ``সময়ের শেষে''
তার দিনলিপি সম্পূর্ণ হবে।
% BN-SRC-815 TARGET-CORRECTION-BEGIN
\par\noindent\textbf{পাঠ-সংশোধন।} এখানে প্রত্যেক দিনের জন্য কোনো সসীম বছরে তার বিবরণ লেখা
হয়—এই আলাদা আলাদা দাবি করা হচ্ছে। এমন একটি সসীম বছর
আছে যখন সব দিনের বিবরণ লেখা শেষ—সেই দাবি নয়।
``সময়ের শেষে'' বলতে সব সসীম বছরে লেখা অংশের সংযোগ বোঝানো
হয়েছে; কোনো শেষ সসীম বছর নয়।
% BN-SRC-815 TARGET-CORRECTION-END

তাহলে এই ভাবনার সঙ্গে কতটা তফাত যে $\kappa$
পর্যন্ত $\alpha$ সর্বদা $\aleph_\alpha$-র চেয়ে ছোট,
এমনকি ক্রমশ আরও অনেক ছোট; অথচ সেই বিন্দুতে এসে
দুই পক্ষ সমান হয় এবং $\kappa=\aleph_\kappa$?

এটি সরিয়ে রেখে, এবং $\ZFC$ মেনে নিয়ে,
স্থিরবিন্দু-নির্মাণের আরও কয়েকটি মজার ফল দিয়ে
শেষ করি। পরের তিনটি ফল স্বজ্ঞায় দেখায়:
স্তরক্রমের কোনো একটি অ-তুচ্ছ পর্যায়ে তার
প্রস্থ ও উচ্চতা সমান হয়।
% BN-SRC-812 TARGET-CORRECTION-BEGIN
\par\noindent\textbf{পাঠ-সংশোধন।} নিচের মূল নির্মাণে শুরু ছিল $\tau_0(A)=\card{A}$। কিন্তু
$\card{A}=\beth_{\card{A}}$ হলে তখন সব পদই $\card{A}$,
তাই ঘোষিত $\card{A}<\tau(A)$ মিথ্যা; পরের $W$-অনুক্রমও
একটি স্থিরবিন্দুতে আটকে যায়। এখানে শুরুটি
$\tau_0(A)=\cardsucc{\card{A}}$ করা হয়েছে।
$\beth_\xi\geq\xi$ এবং বেথ-অপেক্ষকের কঠোর বৃদ্ধিতে
পদগুলি অ-হ্রাসমান। কখনও দুটি পরপর পদ সমান হলে সেই পদই
স্থিরবিন্দু; নইলে পদগুলি সংযোগে সহশেষ, আর বেথের সীমা-অবিচ্ছিন্নতা
দেয় $\beth_{\tau(A)}=\tau(A)$। উভয় ক্ষেত্রেই
$\tau(A)\geq\tau_0(A)>\card{A}$। ফলে সব $W_\alpha$
অঙ্কবাচক সংখ্যা, $W_{\alpha+1}>W_\alpha$, এবং সীমায় সংযোগ
নেওয়ায় $\alpha<\beta$ হলে $W_\alpha<W_\beta$। ধনাত্মক
উত্তরসূরি সূচকে বেথ-স্থিরবিন্দু পাওয়া যায়; ধনাত্মক সীমা সূচকে
আগের ধনাত্মক স্থিরবিন্দুগুলির সহশেষ সংযোগও স্থিরবিন্দু।
$W_0=0$ স্থিরবিন্দু নয়, তবে $\card{V_{W_0}}=W_0$ সরাসরি সত্য।
% BN-SRC-812 TARGET-CORRECTION-END

\begin{prop}\ollabel{bethfixed}
একটি $\beth$-স্থিরবিন্দু আছে, অর্থাৎ এমন
$\kappa$ আছে যার জন্য $\kappa=\beth_\kappa$।
\end{prop}

\begin{proof}
\olref{alephfixed}-এর মতোই; সেখানে ``$\aleph$''-এর
বদলে ``$\beth$'' ব্যবহার করো।
\end{proof}

\begin{prop}\ollabel{stagesize}
$\card{V_{\omega+\alpha}}=\beth_{\alpha}$।
$\omega\ordtimes\omega\leq\alpha$ হলে
$\card{V_\alpha}=\beth_\alpha$।
\end{prop}

\begin{proof}
প্রথম দাবিটি সহজ ট্রান্সফাইনাইট আবেশে পাওয়া যায়।
% BN-SRC-814 TARGET-CORRECTION-BEGIN
\par\noindent\textbf{পাঠ-সংশোধন।} আবেশের ধাপগুলি হলো: $V_\omega$ সসীম পর্যায়গুলির সংযোগ;
প্রত্যেক সসীম পর্যায় সসীম, আর $V_\omega$ সব স্বাভাবিক সংখ্যা
ধারণ করে। তাই চয়নসহ গণনাযোগ্য সংযোগের হিসাবে তার
অঙ্কবাচকতা $\omega=\beth_0$। উত্তরসূরিতে ঘাতসেটের সংজ্ঞাই
ফল দেয়। অশূন্য সীমা $\lambda$-য়
$V_{\omega+\lambda}=\bigcup_{\delta<\lambda}V_{\omega+\delta}$।
ধরি $\mu=\bigcup_{\delta<\lambda}\beth_\delta=\beth_\lambda$।
প্রত্যেক আগের পর্যায়ের অঙ্কবাচকতা $\mu$-র বেশি নয়, এবং
$\beth_\delta\geq\delta$ থেকে $\lambda\leq\mu$, তাই
$\card{\lambda}\leq\mu$। চয়নের সাহায্যে প্রতিটি পর্যায় থেকে
$\mu$-তে একৈক চিত্রণ বেছে নিয়ে সংযোগটিকে
$\lambda\times\mu$-তে পাঠানো যায়: সদস্যটির উপস্থিতির
ক্ষুদ্রতম সূচক প্রথম স্থানাঙ্ক। এই গুণফলের অঙ্কবাচকতা $\mu$
হয়, কারণ $\mu$ অসীম ও $1\leq\card{\lambda}\leq\mu$।
অন্যদিকে সংযোগ প্রতিটি আগের পর্যায় ধারণ করে, তাই তার
অঙ্কবাচকতা সব $\beth_\delta$-র অন্তত সমান, অর্থাৎ অন্তত
$\mu$। দুই দিক মিলে সীমা-ধাপ। শেষ অনুসিদ্ধান্তে বেথ-স্থিরবিন্দু
$\kappa\geq\beth_1>\omega$ একটি অঙ্কবাচক সংখ্যা; তাই
গণনাযোগ্য ক্রমসংখ্যা $\omega\ordtimes\omega$-র চেয়েও বড়।
% BN-SRC-814 TARGET-CORRECTION-END
দ্বিতীয় দাবি আসে, কারণ
$\omega\ordtimes\omega\leq\alpha$ হলে
$\omega+\alpha=\alpha$। এটি দেখাতে
\olref[ord-arithmetic][]{chap}-এর ক্রমসংখ্যার
পাটিগণিতের তথ্য ব্যবহার করি। প্রথমে লক্ষ করো,
$\omega\ordtimes\omega=
\omega\ordtimes(1\ordplus\omega)=
(\omega\ordtimes 1)\ordplus(\omega\ordtimes\omega)=
\omega\ordplus(\omega\ordtimes\omega)$।
এখন $\omega\ordtimes\omega\leq\alpha$,
অর্থাৎ কোনো $\beta$-র জন্য
$\alpha=(\omega\ordtimes\omega)\ordplus\beta$ হলে
$\omega\ordplus\alpha=
\omega\ordplus((\omega\ordtimes\omega)\ordplus\beta)=
(\omega\ordplus(\omega\ordtimes\omega))\ordplus\beta=
(\omega\ordtimes\omega)\ordplus\beta=\alpha$।
\end{proof}

\begin{cor}
কোনো $\kappa$-র জন্য $\card{V_\kappa}=\kappa$।
\end{cor}

\begin{proof}
\olref{bethfixed} অনুযায়ী একটি $\beth$-স্থিরবিন্দু
$\kappa$ নিই। স্পষ্টত $\omega\ordtimes\omega<\kappa$।
তাই \olref{stagesize} থেকে
$\card{V_\kappa}=\beth_\kappa=\kappa$।
\end{proof}

$V_\kappa$-র নিচে যতগুলি পর্যায়,
$V_\kappa$-তে ততগুলি !!{element}s আছে।
অতএব স্বজ্ঞায় $V_\kappa$ যত চওড়া, ততই উঁচু।
এটি ট্রিস্ট্রাম শ্যান্ডির অবস্থার মতো: এক পর্যায়
থেকে পরের পর্যায়ে যেতে \emph{ঘাতসেট} নিই,
তাই প্রতিবার স্তরক্রম \emph{অনেক} বড় হয়।
কিন্তু ``শেষে'', অর্থাৎ $\kappa$ পর্যায়ে,
প্রস্থ উচ্চতার নাগাল পায়।

প্রশ্ন উঠতে পারে: \emph{স্তরক্রমের প্রস্থ ও উচ্চতা
কতবার সমান হয়?} উত্তর: \emph{যতগুলি ক্রমসংখ্যা আছে,
ততবার।} তবে একটু ব্যাখ্যা দরকার।

নিচের মতো একটি পদ $\tau$ সংজ্ঞায়িত করি।
যেকোনো $A$-র জন্য:
\begin{align*}
	\tau_0(A) & \defis \cardsucc{\card{A}}\\
	\tau_{n+1}(A) & \defis \beth_{\tau_n(A)}\\
	\tau(A) & \defis \bigcup_{n < \omega}\tau_n(A)
\intertext{\olref{bethfixed}-এর মতোই, যেকোনো $A$-র জন্য
$\tau(A)$ একটি $\beth$-স্থিরবিন্দু; আর স্পষ্টত
$\card{A}<\tau(A)$। এবার এই পুনরাবৃত্তিমূলক
সংজ্ঞাটি দেখি:}
	W_0 &\defis 0\\
	W_{\alpha + 1} & \defis \tau(W_\alpha)\\
	W_\alpha & \defis \bigcup_{\beta < \alpha} W_\beta
	\text{, $\alpha$ সীমা ক্রমসংখ্যা হলে}
\end{align*}
নির্মাণটি সব ক্রমসংখ্যার জন্য সংজ্ঞায়িত।
% BN-SRC-462: W_0=0 is not a beth fixed point; positive indices are.
স্বজ্ঞায় $W$ ধনাত্মক ক্রমসংখ্যাগুলি থেকে
$\beth$-স্থিরবিন্দুগুলিতে একটি ``!!a{injection}''।
আগের মতোই যেকোনো $\alpha$-র জন্য
$V_{W_\alpha}$ যত চওড়া, তত উঁচু।

\end{document}
